\section{Extreme-scale coordinates for bubble and overlap integrals}

These transformations keep the finite recursive BPH kernels already proved
above, including the bubble's finite $-r/6$ term. They do not substitute a
massless finite part or an MS amplitude. Analytic tail estimates below
control only discarded infinite intervals, separately from quadrature error.

\subsection{A logarithmic radial coordinate for bubble insertions}

Scale loop momentum by $M=m^2$ and use dimensionless $r=k^2/M$. Put
\[
 x=\log(1+r),\quad w=e^{-x},\quad t=1-w,\quad
 \beta(x)=\frac{B_R(r;1)}{1+r}.
\]
The Jacobian is $dr=(1+r)dx$, and $r/(1+r)=t$. The bubble is evaluated
from its same finite closed form or its convergent small-$r$ series. A
large-$x$ form which does not construct $r$ is
\[
 \beta(x)=\frac{(t+4w)^{3/2}}{3\sqrt t}
 \left\{\frac x2+\log\frac{\sqrt t+\sqrt{t+4w}}2\right\}
                  -\frac{4w}{3}-\frac{4t}{9}.
\]
In particular $\beta(x)=x/6-4/9+O((x+1)e^{-x})$.

For the two-point graph let $L=\log(s/M)$ and $q=e^{L-x}$. In terms of
the angular variables $d$ and $\mu=w$ of the preceding bubble chapter,
define
\[
 N=\frac{(1+2w)d^2+6wd+3w}{(1+d)^2}.
\]
Exact substitution into its finite radial formula gives
\[
 \frac{J_{\mathrm{b},2}}s
 =\frac16\int_0^\infty t^2\beta(x)
                  \frac dq\frac{d}{d+w}N\,dx.
\]
This is the formula implemented in the logarithmic route. Ratios rather
than $q$ itself are used when $L-x>0$: with $i_q=e^{x-L}$,
\[
 d/q=\tfrac12\{1-i_q+\sqrt{(1-i_q)^2+4wi_q}\},\quad
 (1+d)^{-1}=\frac{i_q}{d/q+i_q},\quad
 \frac{d}{d+w}=\frac{d/q}{d/q+wi_q}.
\]
For $q\le1$, rationalizing the root gives
$d/q=2w/[\sqrt{(1-q)^2+4wq}+1-q]$. Continuous endpoint limits are used.
No ratio $s/M$ need be formed.

For the three-point insertion let $S=\max(a,b,c)$, $L=\log(S/M)$ and
$(\widehat a,\widehat b,\widehat c)=(a,b,c)/S$. At a Feynman fraction
$u$, put $E=u\widehat a+(1-u)\widehat b$,
$\tau=u(1-u)\widehat c/E$ when $E>0$, and $q=e^{L-x}E$. The exact
angular difference is represented by
\[
 \mathcal E=1-\frac4{(1+q+\sqrt{(1-q)^2+4wq+4tq\tau})^2},\qquad
 J_{\mathrm{b},3}=-\frac12\int_0^\infty t^2\beta(x)
                                \int_0^1\mathcal E\,du\,dx.
\]
The $E=0$ case takes its physical continuous value. Rationalized increments
are used for $q\le1$ and inverse-$q$ expressions for $q>1$. These are
algebraic rewrites of the exact six-dimensional angular average.

The positive panels have width at most four and are split at the physical
layer $x=L$. For $L\ge0$ truncate at $X=L+48$. The established bounds
$F_2\le s^2/(r+M)^3$ and
$0\le R^{-2}-4/(A+W)^2\le2S/R^3$, together with
$B_R(r)\le(r/6)\log(1+r/M)$, give
\[
 \left|\operatorname{tail}\frac{J_{\mathrm{b},2}}s\right|
 \le\frac{X+1}{12}e^{-48},\qquad
 |\operatorname{tail}J_{\mathrm{b},3}|
 \le\frac{X+1}{6}e^{-48}.
\]
These estimates remain meaningful when $S/M$ exceeds binary64 range,
because only $L$ is used.

As an independent two-point asymptotic check, set
\[
 g(x)=\frac16t^2(1+2w)\beta(x),\qquad
 C_2=\int_0^\infty\{g(x)-x/36+2/27\}dx+\frac{29}{144}.
\]
The integrand is regular at zero and exponentially decreasing at infinity.
The transition kernel at fixed $x-L$ is $(1-e^{x-L})^3$ for $x<L$ and
zero for $x>L$. Its correction moments are $-11/6$ and $-85/36$.
Separating the integrable remainder of $g$, or splitting at $x=L/2$,
therefore gives
\[
 \frac{J_{\mathrm{b},2}}s=\frac{L^2}{72}-\frac L8+C_2+o(1),\qquad
 C_2\simeq0.3445763818589743673553120.
\]
The linear and constant transition contributions are respectively
\[
 (1/36)(-11/6),\qquad
 (-2/27)(-11/6)-(1/36)(-85/36)=29/144.
\]
The production calculation integrates the complete massive kernel, rather
than replacing it by this asymptotic check.

\subsection{Exact logarithmic separation for the overlapping graph}

Let $\sigma=s/M$, $L=\log\sigma$. On the original resolved square take
$0<x,w<1$, $u,v\in(0,1)$, and define
\begin{gather*}
 r=1-w,\quad c=1+xr,\quad a=u(1-u),\quad b=v(1-v),\\
 T=a+\frac{r(u-v)^2}{c},\quad
 d_f=r(3+3xr+x^2r^2),\quad i_m=\frac{w}{x+w},\\
 z=\sigma b i_m,\qquad
 \alpha=\frac{\sigma xTi_m}{1+z}.
\end{gather*}
Write $f(z)=(1+z^{-1})\log(1+z)-1=z h(z)$, with its continuous zero
value. The bare-minus-first-counterterm contribution to $J/s$ is exactly
\[
 \int_{[0,1]^4}\frac{2w}{c^3}
 \left[T\{\log(1+z)+f(\alpha)\}-d_f b f(z)\right]dx\,dw\,du\,dv.
\]
The old two sectors are just the two orderings of $(x,w)$ and their
Jacobian; this formula does not change the recursive subtraction.
The common large logarithm has an exactly integrable coefficient. Since
$\int a=\int b=\int(u-v)^2=1/6$, it equals
\[
 \frac13\int_0^1\int_0^1\frac{w}{c^3}
           \{1+r/c-r(c^2+c+1)\}\,dx\,dw=\frac1{36}.
\]
For example integrate $c=1+xr$ first, then the remaining elementary
rational and logarithmic terms. Subtracting $L$ in both logarithmic
factors pointwise and adding $L/36$ analytically prevents an inaccurate
discrete coefficient from being amplified by an enormous $L$.

The remaining negative counterterm, after the trivial second phase
integration, is
\[
 -\int_0^1a\,du\int_0^1\int_0^1
       \frac{2w}{(x+w-xw)^3}
       f\left(\frac{\sigma a xw}{x+w}\right)dx\,dw.
\]
Set $\zeta=1/x+1/w\ge2$. Integrating the other reciprocal variable
gives weight $\zeta(\zeta-2)/(\zeta-1)^3$. Then $v=1/\zeta$ yields
\[
 -\int_0^1a\,du\int_0^{1/2}
       \left\{\frac1v+g(v)\right\}f(\sigma av)dv,
 \quad
 g(v)=\frac1{1-v}+\frac1{(1-v)^2}-\frac1{(1-v)^3}.
\]
The exact primitive of $f(z)/z$ is
$F(z)=-\operatorname{Li}_2(-z)-f(z)$. Define finite corrections
\begin{align*}
 \delta f(z)&=f(z)-\log z+1,\\
 \delta F(z)&=F(z)-\{\tfrac12\log^2z-\log z+1+\pi^2/6\}.
\end{align*}
The counterterm is exactly
\begin{align*}
 &-\frac{L^2}{12}+\frac{19L}{36}+\frac{\pi^2}{72}-\frac{137}{108}\\
 &\quad-\int_0^1a\,\delta F(\sigma a/2)du
       -\int_0^1a\,du\int_0^{1/2}g(v)\delta f(\sigma av)dv.
\end{align*}
Both finite correction integrals are retained at every scale. The constants
follow, in particular, from
\begin{gather*}
 \int a=\tfrac16,\quad \int a\log a=-\tfrac5{18},\quad
 \int a\log^2a=\tfrac{28}{27}-\tfrac{\pi^2}{18},\\
 \int_0^{1/2}g(v)dv=\log2-\tfrac12,\qquad
 \int_0^{1/2}g(v)\log v\,dv=\tfrac12-\tfrac{\pi^2}{12}
                                             -\tfrac12\log^22.
\end{gather*}
Thus the full leading logarithms are
$J/s=-L^2/12+5L/9+O(1)$, but the algorithm does not replace the finite
massive answer by an asymptotic constant.

For stable evaluation of the last correction put $v=e^{-t}/2$ and panelize
$t\ge0$. At $T=L+48$, $z=\sigma av\le e^{-48}/8$. Since $|g|\le2$
and $f(z)\le z/2$, its discarded tail obeys
\[
 |\operatorname{tail}|
 \le\frac{e^{-L-96}}{240}
 +e^{-L-48}\left\{\frac{50+\log2}{6}+\frac5{18}\right\}
 <9e^{-L-48}\qquad(L\ge0).
\]
This is even smaller than a conservative $160e^{-L-48}$ enclosure.
It controls the truncated log-$v$ interval only; phase quadrature and panel
quadrature still require independent refinement. Common Schwinger-regulator
equivalence and the original twelve-dimensional $L^1$ proofs remain those
of the preceding chapters, because every transformation here starts from
their already established finite recursive BPH integrals.

\subsection{A controlled scale-free overlap reference}

An independent reference can integrate the finite constant after taking the
large-scale limit, while bounding the omitted finite-mass terms analytically.
Set $\alpha_\infty=xT/b$, $\ell=\log(bi_m)$ and define the scale-free
four-dimensional constant
\[
 C_\infty=\frac{\pi^2}{72}-\frac{137}{108}
 +\int_{[0,1]^4}\frac{2w}{c^3}
       \{T[\ell+f(\alpha_\infty)]-d_f b(\ell-1)\}
                                        dx\,dw\,du\,dv.
\]
Its logarithmic endpoint terms are integrable. This reference uses a
scale-free constant, whereas the production calculation retains both finite
counterterm corrections and the finite-$z$ raw expression.

For $y>0$, $0\le y f'(y)=1-\log(1+y)/y\le1$. Since
$\alpha=\alpha_\infty z/(1+z)$, this proves
\[
 0\le\log(1+1/z)+f(\alpha)-f(\alpha_\infty)
                                  \le\log(1+1/z).
\]
Thus the absolute difference of the raw density and its limiting
logarithmic polynomial is bounded by
\[
 \frac{2w}{c^3}\{T\log(1+1/z)+d_f b\,\delta f(z)\}.
\]
For $0<\theta<1$ one has
\[
 \log(1+1/z)\le z^{-\theta}/\theta,\qquad
 \delta f(z)\le
       \left(\frac1\theta+\frac1{1-\theta}\right)z^{-\theta}.
\]
Take $\theta=1/4$. The elementary bounds $T\le5/4$, $d_f\le7$,
$2w/c^3\le2w$, and
$z^{-1/4}\le\sigma^{-1/4}(2/(bw))^{1/4}$ give an integrated raw error
at most $K_{\rm raw}\sigma^{-1/4}$, where
\[
 K_{\rm raw}=
 \frac{2^{5/4}}{7/4}
 \left\{5B(3/4,3/4)+\frac{112}{3}B(7/4,7/4)\right\}<32.
\]
Indeed $B(3/4,3/4)\le4\sqrt2/3$ by splitting at $1/2$, and
$B(7/4,7/4)\le1/(2\sqrt2)$ by the maximum of $v(1-v)$.
Hence $K_{\rm raw}\le(64/7)2^{7/4}<32$, without a numerical special-function
enclosure.

For the complete, untruncated counterterm correction, take $\theta=3/4$.
The dilogarithm inversion identity gives
\[
 |\delta F(z)|\le\frac{64}{9}z^{-3/4},\qquad
 \delta f(z)\le\frac{16}{3}z^{-3/4}.
\]
Using $|g(v)|\le2$ on $[0,1/2]$ and integrating $a^{1/4}$ yields the
bound $K_{\rm ct}\sigma^{-3/4}$, with
\[
 K_{\rm ct}=\frac{256}{9}2^{3/4}B(5/4,5/4)<32.
\]
For example Jensen's inequality gives
$B(5/4,5/4)\le6^{-1/4}$, and
$(256/9)(4/3)^{1/4}<32$.
Consequently, for every $\sigma\ge1$,
\[
 \left|\frac{J_{\mathrm{ov}}}{s}
       -\left(-\frac{L^2}{12}+\frac{5L}{9}+C_\infty\right)\right|
 \le32\sigma^{-1/4}+32\sigma^{-3/4}.
\]
At $\sigma=10^{308}$ the bound is less than $3.201\times10^{-76}$.
This controls the whole integrated scale remainder. Numerical quadrature of
$C_\infty$ still has its own refinement evidence and is not made rigorous
by this scale estimate. It can, however, provide a materially different
reference for the production finite-mass formula at extreme ratios.
